\chapter*{Polar Coordinates and Polar Curves}
\addcontentsline{toc}{chapter}{Polar Coordinates and Polar Curves}

\section*{Polar Coordinates}

In polar coordinates, a point is represented by $(r,\theta)$, where $r$ is the directed distance from the origin and $\theta$ is the angle measured from the positive $x$-axis.

The conversion formulas are
$$x=r\cos\theta, \qquad y=r\sin\theta, \qquad r^2=x^2+y^2.$$ 

Polar coordinates are not unique. For example,
$$(r,\theta), \qquad (r,\theta+2\pi), \qquad (-r,\theta+\pi)$$
all describe the same point.

\section*{Graphing Polar Curves}

A polar curve is given by an equation of the form
$$r=f(\theta).$$ 
To sketch it effectively, one usually studies:
\begin{itemize}
    \item symmetry,
    \item intercepts and zeros of $r$,
    \item the maximum and minimum values of $r$,
    \item the behavior over one full period.
\end{itemize}

Common examples include circles, roses, and cardioids such as
$$r=a(1+\sin\theta), \qquad r=a(1+\cos\theta).$$ 

\section*{Area in Polar Coordinates}

When $\alpha\le \theta\le \beta$, the area enclosed by the polar curve $r=f(\theta)$ is
$$A = \frac12\int_{\alpha}^{\beta} r^2\,d\theta.$$ 

\begin{remark}
This formula comes from approximating the region by thin sectors of area
$\frac12 r^2\Delta\theta$.
\end{remark}

\section*{Arc Length in Polar Coordinates}

If $r=f(\theta)$ is differentiable, the arc length from $\theta=\alpha$ to $\theta=\beta$ is
$$L = \int_{\alpha}^{\beta} \sqrt{r^2+\left(\frac{dr}{d\theta}\right)^2}\,d\theta.$$ 

This formula is analogous to the parametric formula because the polar curve can be rewritten parametrically as
$$x=r(\theta)\cos\theta, \qquad y=r(\theta)\sin\theta.$$ 

\section*{Summary of the Unit}

The final part of the course broadens the language of calculus:
\begin{itemize}
    \item parametric equations describe curves using a parameter,
    \item polar coordinates describe points using distance and angle,
    \item derivatives, areas, and lengths all adapt naturally to these new representations.
\end{itemize}